\begin{center}7.\end{center}
\noindent\hspace*{\parindent}Bei der zweiten, der summatorischen, Begründungsart liegt stillschweigend die Voraussetzung in Beziehung auf die Function \(ft\) zum Grunde, dass ihr
\clearpage\SourceLeaf{421}\ArithmeticHead{414}{L}{ANALYSIS. NACHLASS.}
\noindent Werth zwischen den Grenzen \(t=t_0\) und \(t=T\) überall ein bestimmter und zwar ein endlicher sei. Ich füge darüber einige Erläuterungen bei.

Die Forderung, dass der Werth von \(ft\) ein bestimmter sei, ist schon von selbst erfüllt, wenn diese Function in die Classe derjenigen gehört, welche \textsc{Euler}, nicht eben glücklich, mit dem Namen einförmige\textsuperscript{*}) belegt hat, und die man treffender einwerthige nennen könnte, da eben ihr Charakter darin besteht, dass jedem Werthe der Veränderlichen nur Ein Werth der Function entspricht.

Ist hingegen \(ft\) eine vielwerthige Function, so ist erforderlich, dass nach einem Principe bestimmt sei, welcher aus den verschiedenen Werthen der Function für jeden Wert von \(t\) zwischen den Grenzen \(t_0\) und \(T\) gelten soll. Das Princip, dass die Werthe der Function nach der Stetigkeit zusammenhängen sollen, wird, sobald \(ft\) für einen der Werthe von \(t\) gegeben ist, zur Bestimmung des Zuges ausreichen, so lange man nicht auf einen Werth von \(t\) kommt, für welchen der betreffende Werth von \(ft\) mit einem andern zusammenfällt, (also zwei Züge einander kreuzen), in welchem Falle über das ferne Fortschreiten eine anderweitige Bestimmung getroffen werden muss. Übrigens verstehe ich das Fortschreiten nach der Stetigkeit nur so, dass unendlich kleinen Veränderungen von \(t\) unendlich kleine Änderungen von \(ft\) entsprechen, und sehe einen Übergang von reellen Werthen zu imaginären \textls[120]{nicht} als eine Unterbrechung der Stetigkeit an.

Die zweite Bedingung, dass der Werth von \(ft\), so lange \(t\) die Grenzen \(t_0\) und \(T\) nicht überschreitet, überall endlich sei, ist übrigens keine absolut nothwendige. Man nehme an, dass \(ft\) unendlich werde für \(t=t'\), aber für alle zwischen \(t_0\) und \(t'\) liegende Werthe endlich bleibe. Ist \(t'-\omega\) ein solcher Werth, so hat \(\displaystyle\int\limits_{t_0}^{t'-\omega} ft\,{.}\,dt\) einen bestimmten endlichen Werth, und wenn derselbe bei unendlich abnehmendem \(\omega\) nicht über alle Grenzen wächst, sondern sich nur einer endlichen Grenze \(G\) unendlich nähert, so wird es erlaubt
\par\vfill\noindent\rule{.25\linewidth}{.4pt}\par\smallskip
{\footnotesize\noindent\hspace*{\parindent}\textsuperscript{*}) Angemessener würde es sein, einförmige oder einformige Functionen diejenigen zu nennen, die für alle Werthe von \(x\) nach einem einzigen Gesetze gebildet werden; zweiförmige, vielförmige hingegen wären dann die, wo für die Werthe zwischen gewissen Grenzen ein Gesetz gilt, zwischen andern ein anderes gilt, kurz was man sonst functiones discontinuae nannte. In neuester Zeit hat [man] aber die Benennung continuirliche und discontinuirliche Function in einer andern Bedeutung zu nehmen angefangen.\par}
\clearpage\SourceLeaf{422}\ArithmeticHead{415}{R}{CONVERGENZ DER ENTWICKLUNG PERIODISCHER FUNCTIONEN.}
\noindent sein, unter \(\displaystyle\int\limits_{t_0}^{t'}\) eben diese Grenze zu verstehen. Bleibt dann ferner \(ft\) endlich für alle Werthe von \(t\), welche zwischen \(t'\) und \(T\) fallen, letztern eingeschlossen, und ist \(t'+\omega\) ein solcher Werth, so wird \(\displaystyle\int\limits_{t'+\omega}^{T} ft\,{.}\,dt\) einen bestimmten endlichen Werth haben, und wenn dieser bei unendlich abnehmendem \(\omega\) sich nur einer endlichen Grenze unendlich nähert, so wird diese Grenze \(G'\) als Werth von \(\displaystyle\int\limits_{t'}^{T} ft\,{.}\,dt\) betrachtet werden dürfen, und mit gleichem Rechte die Summe \(G+G'\) als Werth von \(\displaystyle\int\limits_{t_0}^{T} ft\,{.}\,dt\)[\textsuperscript{*})].

Wenn hingegen diese Bedingungen nicht beide zutreffen, d. i. wenn indem \(ft'\) unendlich wird,
\[
\int\limits_{t_0}^{t'-\omega} ft\,{.}\,dt\ \text{und}\ \int\limits_{t'+\omega}^{T} ft\,{.}\,dt,
\]
\noindent bei unendlich abnehmendem \(\omega\), nicht beide endliche Grenzwerthe haben, so bildet in der zweiten Begründungsart der Wert \(t=t'\) eine unübersteigliche Scheidewand, so lange man (wie man bisher immer gethan hat) sich dabei ausschliesslich auf reelle Werthe von \(t\) beschränkt. Es findet sich also hier eine Lücke in der Integralrechnung: denn in allen den unzähligen Fällen, wo \(ft\,{.}\,dt\) nicht als das Differential einer analytischen[\textsuperscript{**})] Function \(Ft\) dargestellt werden kann, d. i. einer solchen, die eine allgemein gültige Werth-Bestimmung für jeden Werth von \(t\) durch algebraische oder ganz eingebürgerte transscendente Operationen darbietet, und wo man also die Bedeutung von \(\int ft\,{.}\,dt\) nur in unserer zweiten Art auffassen kann, bleiben, unter den vorhin bezeichneten Umständen die beiden Parthien der Functionalwerthe, diesseits und jenseits der Scheidewand, ohne Verknüpfung. Folgeweise hat man dann in solchen Fällen öfters unberechtigte Willkür eintreten lassen, während die strenge Wissenschaft für jede analytische Function, die sie in ihren Bereich zieht, eine allgemeingültige Genesis aus einem einheitlichen Princip verlangt. Allerdings ist gegen diese Forderung in der neuern Mathematik fort und fort vielfach gefehlt.
\par\vfill\noindent\rule{.25\linewidth}{.4pt}\par\smallskip
{\footnotesize
\noindent\hspace*{\parindent}[\textsuperscript{*}) In der Handschrift lautet die untere Grenze des Integrals \(t'\); im Vorhergehenden schrieb \textsc{Gauss} für die untere Grenze abwechselnd bald \(t_0\), bald \(t^0\).]\par
\noindent\hspace*{\parindent}[\textsuperscript{**}) Früher hieß es »entwickelten«, was \textsc{Gauss} durchgestrichen hat.]\par}
